\subsection{Condiciones experimentales}

Los experimentos se ejecutaron en el equipo descrito en la Sección~3.3, con un
límite de 60 segundos por resolución del modelo entero y una pasada por
instancia, dado el carácter determinista del método. Para cada instancia se
registró el valor de la solución obtenida ($z^{*}$), el número de iteraciones
del ciclo de generación de cortes, el número total de desigualdades RCI
incorporadas, el tiempo total de resolución y la brecha relativa (gap) respecto
del valor óptimo de referencia leído desde el archivo \texttt{.sol} de CVRPLIB.

\subsection{Resultados}

La Tabla~\ref{tab:resultados} resume los resultados sobre las cinco instancias
evaluadas.

\begin{table}[H]
\centering
\caption{Resultados computacionales sobre instancias de CVRPLIB (Sets E y A).
Límite de 60~s por resolución del modelo entero.}
\label{tab:resultados}
\begin{tabular}{lrrrrrrrl}
\hline
Instancia & $n$ & $m$ & Óptimo & $z^{*}$ & Gap (\%) & Iter. & Cortes & Tiempo (s) \\
\hline
E-n13-k4 & 13 & 4 & 247 & 247 & 0{,}00 & 15  & 47     & 0{,}9 \\
E-n22-k4 & 22 & 4 & 375 & 375 & 0{,}00 & 13  & 84     & 0{,}8 \\
E-n23-k3 & 23 & 3 & 569 & 569 & 0{,}00 & 16  & 135    & 1{,}2 \\
A-n32-k5 & 32 & 5 & 784 & 784 & 0{,}00 & 99  & 788    & 576{,}1 \\
A-n33-k5 & 33 & 5 & 661 & 667 & 0{,}91 & 384 & 2\,846 & 17\,096{,}4 \\
\hline
\end{tabular}
\end{table}

En cuatro de las cinco instancias el método alcanzó el valor óptimo de
referencia con gap nulo. En estos casos, cada resolución del modelo entero
concluyó dentro del límite de tiempo con certificado de optimalidad del solver,
por lo que el valor final corresponde al óptimo global certificado del CVRP.
La instancia \texttt{A-n33-k5} constituye la excepción y se analiza por
separado en la Sección~\ref{sec:limite}.

La Figura~\ref{fig:rutas32} muestra la solución óptima obtenida para
\texttt{A-n32-k5}: cinco rutas que parten y retornan al depósito, cada una con
carga acumulada dentro de la capacidad $D=100$.

\begin{figure}[H]
    \centering
    \includegraphics[width=0.78\textwidth]{figuras/A-n32-k5_rutas_20260611_213833.png}
    \caption{Rutas óptimas obtenidas para la instancia \texttt{A-n32-k5}
    ($z^{*}=784$). El cuadrado indica el depósito. Fuente: elaboración propia.}
    \label{fig:rutas32}
\end{figure}

\subsection{Análisis de la convergencia}

El comportamiento del método sobre \texttt{A-n32-k5} ilustra la mecánica del
esquema de generación iterativa de cortes. El modelo inicial, que contiene
únicamente las restricciones de grado, entrega un valor de 585: esta es la
cota inferior provista por la relajación, un 25{,}4\,\% por debajo del óptimo.
A medida que el procedimiento detecta subtours ilegales y excesos de capacidad
e incorpora las desigualdades RCI correspondientes (Ec.~5), la cota asciende
de forma monótona: 743 tras 10 iteraciones, 780 tras 50, hasta alcanzar 784 en
la iteración 91 y certificar, en la iteración 99, que ningún subconjunto
adicional incumple las restricciones. En total se incorporaron 788
desigualdades de una familia cuyo tamaño completo es exponencial, lo que
confirma la pertinencia de generarlas únicamente cuando la estructura de la
solución lo exige.

El análisis del tiempo por resolución revela dónde reside el costo del
esquema: las primeras iteraciones se resuelven en centésimas de segundo,
mientras que las últimas requieren más de 16 segundos cada una, pues el modelo
acumula cientos de cortes y debe resolverse desde cero en cada iteración. El
esfuerzo computacional, por tanto, no está dominado por la separación de
cortes (de costo polinomial mediante componentes conexas) sino por las
re-resoluciones sucesivas del modelo entero.

Las instancias del Set E exhiben el mismo patrón en miniatura: convergen en
13--16 iteraciones y alrededor de un segundo, con la cota ascendiendo
monótonamente hasta el óptimo. La comparación entre ambos grupos anticipa el
crecimiento abrupto del esfuerzo con el tamaño de la instancia.

\subsection{El límite del esquema: la instancia A-n33-k5}
\label{sec:limite}

Con un nodo más que \texttt{A-n32-k5}, la instancia \texttt{A-n33-k5} no pudo
resolverse a optimalidad dentro del presupuesto experimental. La traza de
ejecución muestra que, a partir de la iteración 148 aproximadamente, cada
resolución del modelo entero agotó el límite de 60 segundos sin que el solver
certificara la optimalidad del subproblema, retornando en su lugar la mejor
solución entera encontrada. La evidencia de esta pérdida de certificación es
visible en la propia traza: hasta ese punto el valor de la cota asciende de
forma monótona (496 a 647), comportamiento propio de óptimos de modelos
progresivamente más restringidos, mientras que después oscila de manera
errática (entre 650 y 670), lo que corresponde a soluciones factibles no
necesariamente óptimas de cada subproblema.

El procedimiento concluyó tras 384 iteraciones y 4{,}7 horas con una solución
factible de costo 667, un 0{,}91\,\% sobre el óptimo de referencia (661). Cabe
subrayar la distinción: se trata de una solución factible de buena calidad,
pero sin certificado de optimalidad, pues la última resolución no concluyó
dentro del límite de tiempo. Este resultado delimita empíricamente, para la
configuración utilizada, la frontera de instancias resolubles a optimalidad
por el esquema clásico: el salto de $n=32$ a $n=33$ multiplicó el tiempo de
resolución por un factor cercano a 30 sin alcanzar el certificado.
